\section{Research Background}

In 2015, the United Nations adopted the 2030 Agenda for Sustainable Development. Among the 17 Sustainable Development Goals (SDGs), SDG 7 aims to ensure access to affordable, reliable, sustainable, and modern energy for all. Buildings account for a significant share of global energy demand, making them a critical sector for achieving SDG7. However, major challenges persist, including excessive building energy consumption and rapid growth of air-conditioning demand, unequal access to clean energy, and the lack of optimisation strategies at the community scale, which hinders the balance between energy conservation and thermal comfort. Despite international efforts, progress in improving energy efficiency remains significantly insufficient, as highlighted in \textit{The Sustainable Development Goals Report 2025}. This underscores the urgent need for integrated research examining the interactions among urban microclimate, building energy performance, and human comfort.

Current global research addressing urban energy optimisation primarily focuses on simulation and field measurements to improve thermal comfort and building energy performance at the community level. Previous studies have shown that community design, vegetation layout, and building form can effectively improve the thermal environment and reduce energy demand \citep{oke1981canyon,wang2025impact,zheng2025reanalysis}. Nevertheless, most existing studies remain limited to single-scale spatial analysis and lack integrated approaches that simultaneously link neighbourhood-level thermal comfort with building-level energy consumption. This gap hinders the formulation of holistic urban planning strategies capable of addressing both energy savings and environmental quality.

\subsection{Theories and Models of Thermal Comfort}

As research on thermal comfort has expanded, its theories and models have evolved from focusing on steady-state, uniform environments to dynamic and non-uniform conditions. \citet{deDear2002thermal}, through the analysis of 21,000 field datasets worldwide, proposed the Adaptive Model, which incorporated outdoor climate and human behavioural and psychological factors into the thermal comfort framework. This model also contributed to the revision of the ASHRAE 55 standard, expanding its applicability to naturally ventilated buildings.

Later, \citet{enescu2017review} systematically reviewed indoor thermal comfort models, highlighting classic models such as Predicted Mean Vote (PMV) and Predicted Percentage of Dissatisfied (PPD), along with their extended versions. The study also pointed out the potential of artificial intelligence technologies in improving thermal comfort control in buildings. \citet{binarti2020review} focused on hot and humid regions, reviewing outdoor thermal comfort indices such as Physiological Equivalent Temperature (PET), Standard Effective Temperature (SET*), and Universal Thermal Climate Index (UTCI). They suggested adjusting the neutral temperature range to better reflect local climate characteristics. \citet{zhao2021thermal} further summarised the development of models for special populations, extreme environments, and data-driven approaches, emphasising that future thermal comfort models should aim for cross-scenario adaptability and real-time control.

Although existing research has broadened the range and applicability of thermal comfort models, challenges remain, including lack of standardisation and insufficient consideration of psychological and behavioural factors.

\subsection{Neighbourhood-scale Microclimate Studies}

Building upon thermal comfort theories, research at the neighbourhood scale has examined how microclimates affect human comfort and building performance. The microclimate of a neighbourhood can directly influence the living environment of people and thereby determine their comfort and health levels. Neighbourhood-scale microclimate studies mainly rely on simulation-based analyses using detailed urban morphology models.

\citet{oke1981canyon}, in a classic study combining physical scale models with field observations, revealed the key role of urban canyon geometry in night-time heat island formation. Oke demonstrated that the urban heat island effect is driven by the accumulation and amplification of heat in urban canyon microclimates. He introduced the Sky View Factor (SVF) to quantify the heat dissipation potential of building clusters, providing a fundamental theoretical basis for later research.

Building on this, \citet{wang2025impact} simulated 25 typical residential communities in Wuhan and found that SVF and green coverage ratio were the most important factors influencing outdoor thermal comfort in both summer and winter. Meanwhile, shape factor (SF) and building density (BD) had significant effects on heating and cooling energy consumption. To address the problem of insufficient urban meteorological data, \citet{zheng2025reanalysis} proposed a method that integrates ERA5 reanalysis data with urban land surface models. By combining UT\&C and Urban Weather Generator (UWG) models, they corrected global datasets to generate high-resolution local urban climate data for building energy simulations.

Overall, existing studies have begun to establish connections among urban geometry, local meteorology, and building energy consumption. However, limitations remain, such as restricted model applicability and insufficient consideration of psychological and behavioural factors. These findings are increasingly relevant for building energy simulation studies, as local climate strongly affects indoor loads. This highlights the importance of integrating neighbourhood climate data into Building Energy Modelling (BEM).

\subsection{Building Energy Simulation}

Research on building energy simulation (BES) has closely followed the development and application of simulation tools. \citet{crawley2001energyplus} introduced EnergyPlus as a new-generation building energy simulation tool, highlighting its advantages in simultaneous load and system calculations, modular architecture, and open interfaces. These features laid a solid technical foundation for the later development of Urban Building Energy Modelling (UBEM) platforms.

\citet{pan2023building} systematically reviewed studies from the past decade and classified the applications of BEM into five main areas: performance-driven design, model-based operational optimisation, digital twins, urban energy modelling, and building-to-grid (B2G) interaction. They reviewed commonly used tools, technologies, and case studies in each area, emphasising the critical role of BEM in scaling applications from individual buildings to neighbourhoods and even entire cities.

\citet{salvalai2024from} used bibliometric analysis to review the evolution of UBEM, comparing nine mainstream platforms, including URBANopt, CityBES, and CitySim. Their findings showed that URBANopt excels in net-zero community and distributed energy network simulations, while CityBES is more effective for retrofit assessments of existing building stocks. However, they also pointed out that UBEM still faces significant challenges, such as insufficient geometric and physical parameter data, high computational demands, lack of cross-platform standards, and difficulties in modelling human behaviour.

\subsection{Coupling Urban Microclimate and Building Energy Models}

Differences between urban climates and the suburban locations where meteorological data are typically collected can significantly affect the accuracy of building energy simulations. To address this issue, researchers in recent years have proposed coupling strategies that integrate urban microclimate characteristics into energy models. \citet{sezer2023urban}, after conducting a systematic literature review, found that compared with standalone simulations, coupled models can more accurately assess the impact of the urban heat island effect on building energy consumption and are particularly suitable for high-density urban areas, where heat accumulation and energy use are strongly interlinked. However, most current studies still focus on one-way coupling, which struggles to capture temporal variability of summer conditions. In contrast, bidirectional coupling can improve accuracy but requires significantly higher computational resources and remains in the exploratory stage.

\subsection{Research Gap}

Despite significant progress in thermal comfort research, neighbourhood-scale microclimate studies, and building energy simulations, existing work remains largely fragmented. Most studies focus on single-scale analysis and adopt one-way coupling. This approach fails to capture the dynamic feedback between outdoor thermal environments and building energy performance. Moreover, human behavioural and psychological factors, which are critical determinants of thermal comfort, are often overlooked. To fill this research gap, it is essential to develop an external bidirectional coupling platform that integrates urban microclimate models (UCMs) and building energy models (BEMs), while explicitly incorporating human factors, to provide a more comprehensive and realistic assessment of energy--comfort interactions across scales.
